% Minden, ami százalék jel után áll, az megjegyzés (nem jelenik meg a pdf-ben).
% Az első néhány sor tex-nikai dologkat tartalmaz, nem kell foglalkozni vele.
\documentclass[a4paper]{article}
\usepackage[T1]{fontenc}
\usepackage[utf8]{inputenc}
\usepackage{amsfonts}
\usepackage{amssymb}
\usepackage{amsmath}

\begin{document}

% Itt kezdődik a tényleges szöveg.

Minden matematikai jelölést dollárjelek közé kell tenni, 
még akkor is, ha csak egyetlen betűről van szó. 
Hasonlítsuk össze a következő két mondatot:\\ % új sort kezdünk
Jelölje a az A csúccsal szemközti oldalt.\\
Jelölje $a$ az $A$ csúccsal szemközti oldalt.

% ha egy üres sort hagyunk, akkor új bekezdés kezdődik
A LaTeX parancsok backslash karakterrel kezdődnek, pl. $\oplus$, $\rightarrow$, $\sin$, $\log$, $\forall$, $\exists$, $\infty$, $\overline{U \pm V}$ stb. A kapcsos zárójelek a programnyelv részét képezik, ezekkel lehet tagolni az inputot (lásd a lenti példákat). Ha ismeretlen parancsot írunk be (vagy elgépelünk egy ismert parancsot) akkor hibajelzést kapunk, éppúgy, mint akkor, ha a kapcsos zárójelek vagy a dollárjelek nincsenek megfelelően párosítva, vagy ha egy matematikai jelölést nem teszünk dollárjelek közé.

\bigskip % kihagyunk egy kis helyet

Néhány gyakran használt jelölés:
\begin{itemize} % felsorolás kezdete

\item kitevő: % minden pontot \item-mel kezdünk
$x^2+a^23+a^{23}+b^{p^2}$;

\item index: 
$x_2+a_23+a_{23}+b_{p_2}$;

\item kitevő és index vegyesen: 
$x_3^2+x^2_3+a^{b_c}+a_{b^c}$;

\item műveleti jelek: 
$a+b, a-b, a \cdot b, a/b, f \circ g, A \cup B, A \cap B, A \setminus B$;

\item tört: 
$\frac{a+bi}{c+di}=\frac{(ac-bd)+(ad+bc)i}{c^2+d^2}$;

\item gyökvonás:
$\sqrt{-1}+\sqrt[3]{2}+\sqrt[n+m]{a+b+c}$

\item relációs jelek: 
$2 \leq 3, x \geq y, a 
eq b, \emptyset \subseteq H, B \supseteq A, A \subset B, A \supset B$;
\item \emph{és így tovább} jele (három pont): 
$a_1+ \ldots +a_n$,
$a_1+ \cdots +a_n$,
$a_1 \cdot\ldots\cdot a_n$,
$a_1 \cdot\cdots\cdot a_n$ (ez utóbbit inkább kerüljük!);

\item kongruencia: 
$a \equiv b \pmod m \iff m \mid a-b$;

\item halmaz megadása: $\{ 2k \mid k \in \mathbf{Z} \}$ vagy $\{ n : n \text{ páros egész szám} \}$;

\item görög kisbetűk:
$\alpha,\beta,\gamma,\delta,\epsilon,\zeta,\eta,\theta,\iota,\kappa,\lambda,\mu,
u,\xi,o,\pi,\rho,\sigma,\tau,\upsilon,\phi,\chi,\psi,\omega$ 
továbbá $\varepsilon,\vartheta,\varkappa,\varrho,\varphi$;

\item görög nagybetűk:
$A,B,\Gamma,\Delta,E,Z,H,\Theta,I,K,\Lambda,M,N,\Xi,O,\Pi,P,\Sigma,T,\Upsilon,\Phi,X,\Psi,\Omega$.

\end{itemize} % felsorolás vége

\bigskip

Ékezetes betűk:

árvíztűrő tükörfúrógép

ÁRVÍZTŰRŐ TÜKÖRFÚRÓGÉP

Ha nagyon nem akarnak működni az ékezetes betűk, végső esetben LaTeX-parancsokkal is megkaphatók:

\'{a}rv\'{\i}zt\H{u}r\H{o} t\"{u}k\"{o}rf\'{u}r\'{o}g\'{e}p

\'{A}RV\'{I}ZT\H{U}R\H{O} T\"{U}K\"{O}RF\'{U}R\'{O}G\'{E}P

\end{document}